\chapter*{Lời Mở Đầu}
\addcontentsline{toc}{chapter}{Lời Mở Đầu}
\markboth{Lời Mở Đầu}{Lời Mở Đầu}

\begin{truyen}{Lời Mở Đầu}{Opening Words}
\chuhoa{C}{ó những câu học sinh nói ra nghe qua tưởng chỉ là cãi cùn. Nhưng nghe kỹ hơn, ta thấy bên dưới thường là một mớ rất thật: áp lực điểm số, nỗi sợ bị coi thường, thói quen trì hoãn, sự mệt mỏi, hoặc đơn giản là không biết nói cho đúng ý mình.}
\textit{(Some student lines sound like mere stubborn excuses. Listen more carefully and you often hear something real underneath: pressure, fear of disrespect, procrastination, exhaustion, or simply poor expression.)}

Vì vậy bản V4 này không đi theo một giọng duy nhất. Nó giữ nguyên mùi lớp học, nhưng mở rộng góc nhìn: có học trò ngụy biện, có giáo viên mắc kẹt trong quy định, có phụ huynh đẩy áp lực đi vòng, và có những ảo tưởng tương lai nghe rất đẹp nhưng đứng không vững. Chính ở những va chạm đó, người dạy lẫn người học mới nhìn rõ mình hơn.
\textit{(That is why this V4 does not follow a single voice. It keeps the smell of the classroom while widening the lens: students making excuses, teachers trapped in rigid rules, parents redirecting pressure, and future dreams that sound beautiful but stand on weak ground. Those collisions are where both teachers and students see themselves more clearly.)}

Mỗi đoạn trong cuốn này vẫn ngắn. Không dài dòng, không lên lớp quá tay. Chỉ là một tình huống, một câu đáp, một cú chững lại, rồi rút ra điều đáng nghĩ. Vì ngoài đời, rất nhiều bài học bắt đầu từ một câu nói tưởng nhỏ mà để lộ cả một kiểu nghĩ phía sau.
\textit{(Each piece in this book is still brief. No long sermons, no over-explaining. Just one situation, one reply, one pause, and then the point worth reflecting on. In real life, many lessons begin with a small sentence that reveals a whole mindset behind it.)}
\end{truyen}

\begin{baihoc}
Khi lý sự bị bóc trần, điều còn lại không phải là thắng thua, mà là cơ hội để lớn lên.
\textit{(When an excuse is exposed, what remains is not victory or defeat, but a chance to grow.)}
\end{baihoc}
\ngancach